\section{Final assembly and the quantitative constant}
\label{sec:final-assembly-v3}

We assemble the ordinary coloring lower bound, the signed cocoloring upper
bound, and the exact constant ledger. No independence between the two final
events is required.

Let $r_+(n)$ denote the ordinary first-moment root and let
$r_4^{\mathrm{co}}(n)$ denote the signed four-size root. Section~5 proves
\begin{equation}
  r_+(n)-r_4^{\mathrm{co}}(n)
  =
  \left[
    \frac{(\log2)^2}{4}A_4(\delta_n)+o(1)
  \right]
  \frac{n}{(\log n)^3},
  \qquad
  A_4(\delta):=\log2-D_4(\delta).
  \label{eq:phase-resolved-root-gap-v3}
\end{equation}
The tangent-rounded signed profile is placed at
\[
  k_{\mathrm{co}}
  =
  \left\lceil
    \frac{r_4^{\mathrm{co}}(n)+r_+(n)}2
  \right\rceil+O(1),
\]
where the bounded correction enforces the two exact profile conservation laws.
Consequently
\begin{equation}
  r_+(n)-k_{\mathrm{co}}
  =
  \left[
    \frac{(\log2)^2}{8}A_4(\delta_n)+o(1)
  \right]
  \frac{n}{(\log n)^3}.
  \label{eq:midpoint-retained-gap-v3}
\end{equation}
Indeed, taking the midpoint divides the leading root separation by two, while
the ceiling and tangent correction contribute only $O(1)$ classes.

Section~4 gives a deterministic integer $k_\chi^-$ such that
\[
  \mathbb P\bigl(\chi(G_n)>k_\chi^-\bigr)\longrightarrow1
\]
and
\begin{equation}
  k_\chi^- - k_{\mathrm{co}}
  =
  \left[
    \frac{(\log2)^2}{8}A_4(\delta_n)+o(1)
  \right]
  \frac{n}{(\log n)^3}.
  \label{eq:integer-location-gap-v3}
\end{equation}

Proposition~\ref{prop:normalized-second-moment-v3} and Paley--Zygmund give a
signed witness with probability at least $e^{-\Lambda_n}$, where
\[
  \Lambda_n=o\!\left(\frac{n}{(\log n)^4}\right).
\]
Section~10 applies the rare-seed amplifier and the simultaneous leftover
coloring theorem. It produces a deterministic $a_n$ such that
\[
  a_n=o\!\left(\frac{n}{(\log n)^3}\right)
\]
and
\[
  \mathbb P\bigl(\zeta(G_n)\le k_{\mathrm{co}}+a_n\bigr)
  \longrightarrow1.
\]
Intersecting the two high-probability events by a union bound and using
\eqref{eq:integer-location-gap-v3} proves the phase-resolved conclusion
\begin{equation}
  \chi(G_n)-\zeta(G_n)
  \ge
  \left[
    \frac{(\log2)^2}{8}A_4(\delta_n)-o(1)
  \right]
  \frac{n}{(\log n)^3}
  \label{eq:phase-resolved-final-v3}
\end{equation}
with high probability.

\subsection{Exact four-support certificate}

For completeness, we record the elementary uniform certificate used to turn
\eqref{eq:phase-resolved-final-v3} into a fixed constant. Let $\lambda_4$ be
the tilt of the limiting four-deficit optimizer. Exact rational estimates give
\[
  \frac{49}{20}\log2
  <
  \lambda_4
  <
  \frac{83}{20}\log2.
\]
Split the omitted limiting mass at $(29/10)\log2$. If $L(\lambda)$ and
$H(\lambda)$ denote the omitted low- and high-deficit ratios, respectively,
the following four bounds hold:
\[
  L\!\left(\frac{49}{20}\log2\right)<\frac{263}{1000},
  \qquad
  H\!\left(\frac{29}{10}\log2\right)<\frac3{200},
\]
\[
  L\!\left(\frac{29}{10}\log2\right)<\frac{33}{250},
  \qquad
  H\!\left(\frac{83}{20}\log2\right)<\frac{29}{200}.
\]
Monotonicity of the two omitted tails therefore gives
\[
  L(\lambda_4)+H(\lambda_4)<\frac{139}{500}.
\]
Evaluating the unrestricted dual objective at the four-support optimizer yields
\[
  D_4(\delta)
  \le
  \log\bigl(1+L(\lambda_4)+H(\lambda_4)\bigr)
  <
  \log\!\left(\frac{639}{500}\right).
\]
Hence, uniformly for $0\le\delta\le1$,
\begin{equation}
  A_4(\delta)
  =
  \log2-D_4(\delta)
  >
  \log\!\left(\frac{1000}{639}\right).
  \label{eq:exact-four-support-certificate-v3}
\end{equation}

Combining \eqref{eq:phase-resolved-final-v3} and
\eqref{eq:exact-four-support-certificate-v3} gives
\[
  \chi(G_n)-\zeta(G_n)
  \ge
  \left[
    \frac{(\log2)^2}{8}
    \log\!\left(\frac{1000}{639}\right)-o(1)
  \right]
  \frac{n}{(\log n)^3}
\]
with high probability. The fixed coefficient is
\[
  \frac{(\log2)^2}{8}
  \log\!\left(\frac{1000}{639}\right)
  =0.026896409808379\ldots.
\]

\ifErdosProofClosed
This proves Theorem~\ref{thm:main-v3}. Since $n/(\log n)^3\to\infty$, the
chromatic--cochromatic difference tends to infinity with high probability
along the full sequence of integers.
\else
Equations \eqref{eq:phase-resolved-final-v3} and
\eqref{eq:exact-four-support-certificate-v3} prove the target theorem
conditional on the four formal gates stated in Appendix~A. The publication
switch remains disabled until those gates are independently validated.
\fi

\begin{corollary}[Simultaneous complement form]
Under the hypotheses of Theorem~\ref{thm:main-v3},
\[
  \min\{\chi(G_n),\chi(\overline{G_n})\}-\zeta(G_n)
  \ge
  \left[
    \frac{(\log2)^2}{8}
    \log\!\left(\frac{1000}{639}\right)-o(1)
  \right]
  \frac{n}{(\log n)^3}
\]
with high probability.
\end{corollary}

\begin{proof}
The law of $\overline{G_n}$ is again $G(n,1/2)$, and
$\zeta(\overline G)=\zeta(G)$. Apply the theorem simultaneously to $G_n$ and
its complement and intersect the two events by a union bound.
\end{proof}
